\textbf{Proof.}
On \(\mathcal E_{1n}\cap\mathcal E_{2n}\), the Bahadur representation, covariance-matched Gaussian and multiplier approximations, anti-concentration, and ratio-consistent studentization imply
\[
 \liminf_{n\to\infty}\inf_QP_Q\left\{
 |\widehat\tau_n(v)-\Theta_{v,h_n}|
 \le\frac{\widehat c_{1-\alpha}\widehat s_n(v)}{\sqrt{n_2h_n}}
 \text{ for every }v\in\mathcal I\right\}\ge1-\alpha. \tag{1}
\]
The center, score, and estimator use the same reference-pilot-selected nonzero winding.  The bias lemma gives
\[
 \sup_{v\in\mathcal I}|\Theta_{v,h_n}-\tau_0(v)|
 \le B_{\mathrm{bias}}L_\tau h_n^\beta. \tag{2}
\]
The triangle inequality proves coverage on success.  On either failure the procedure returns \([-C_\tau,C_\tau]\), which covers by the slope-range restriction.  This proves uniform asymptotic honesty.

On success the sup-width is at most
\[
 \frac{2\widehat c_{1-\alpha}\sup_v\widehat s_n(v)}{\sqrt{n_2h_n}}
 +2B_{\mathrm{bias}}L_\tau h_n^\beta. \tag{3}
\]
On failure it is \(2C_\tau\).  The Bahadur lemma therefore gives
\[
 \sup_QW_n(\mathcal B,Q)
 \le C\sqrt{\frac{\log n}{n_2h_n}}
 +2B_{\mathrm{bias}}L_\tau h_n^\beta+o(r_n). \tag{4}
\]
For the displayed bandwidth, both leading terms are of order \(r_n\).  The direct conditional-\(L^1\) class is included in this result by \(\text{lem:modifier-shift-attainment-and-inclusion}\), and its constant-shift witness proves conditional strictness under the compatibility stated there.  The construction uses the fixed external code and no outcome regression. \(\square\)